% Part: first-order-logic
% Chapter: axiomatic-deduction
% Section: deduction-theorem-quantifiers

\documentclass[../../../include/open-logic-section]{subfiles}

\begin{document}

\olfileid{fol}{axd}{ddq}

\olsection{પરિમાણકો સાથે નિગમન પ્રમેય}

\begin{thm}[નિગમન પ્રમેય]
\ollabel{thm:deduction-thm-q} જો $\Gamma \cup \{!A\} \Proves !B$,
તો $\Gamma \Proves !A \lif !B$.
\end{thm}

\begin{proof}
ફરી આપણે~$\Gamma \cup \{!A\}$માંથી $!B$ની !!{derivation}ની લંબાઈ પર
ગણિતીય અનુમાન કરીએ છીએ.

અનુમાનના આધાર-કિસ્સાની સાબિતી~\olref[ded]{thm:deduction-thm}ની
સાબિતીમાંની સાબિતી જેવી જ છે.

અનુમાનના પગલા માટે ફરી ધારો કે~$\Gamma \cup \{!A\}$માંથી $!B$ની
!!{derivation}નું છેલ્લું પગલું~$!B$ કોઈ અનુમાન-નિયમ વડે પ્રમાણિત થાય
છે. અનુમાન-નિયમ મોડસ પોનેન્સ હોય તો આપણે
\olref[ded]{thm:deduction-thm}ની સાબિતીની જેમ આગળ વધીએ. અનુમાન-નિયમ
\QR{} હોય તો $!B \ident !C \lif \lforall[x][!D(x)]$ છે અને
!!{derivation}માં અગાઉ $!C \lif !D(a)$ સ્વરૂપનું !!a{formula} આવે છે,
જ્યાં $a$ એ~$!C$, $!A$ કે~$\Gamma$માં આવતું નથી. તેથી
\begin{align*}
  \Gamma \cup \{!A\} & \Proves !C \lif !D(a),\\
  \intertext{અને અનુમાનની પૂર્વધારણા લાગુ પડે છે, એટલે આપણી પાસે}
    \Gamma & \Proves !A \lif (!C \lif !D(a)).\\
  \intertext{નીચેના પરિણામથી}
  & \Proves (!A \lif (!C \lif !D(a))) \lif ((!A \land !C) \lif !D(a))\\
  \intertext{અને મોડસ પોનેન્સ વડે મળે છે}
  \Gamma & \Proves (!A \land !C) \lif !D(a).\\
  \intertext{આઇગનચલ-શરત હજુ લાગુ પડતી હોવાથી, આ !!{derivation}માં \QR{} વડે પ્રમાણિત પગલું ઉમેરીને મળે છે}
    \Gamma & \Proves (!A \land !C) \lif \lforall[x][!D(x)].\\
    \intertext{વળી, આપણી પાસે}
    & \Proves ((!A \land !C) \lif \lforall[x][!D(x)]) \lif (!A \lif (!C \lif \lforall[x][!D(x)])),\\
    \intertext{તેથી મોડસ પોનેન્સ વડે}
    \Gamma & \Proves !A \lif (!C \lif \lforall[x][!D(x)]),
\end{align*}
એટલે $\Gamma \Proves !A \lif !B$.

$!B$ને \QR{} નિયમ વડે પ્રમાણિત કરવામાં આવે પરંતુ તેનું સ્વરૂપ
$\lexists[x][!D(x)] \lif !C$ હોય તે કિસ્સો પ્રશ્ન તરીકે રાખીએ છીએ.
\end{proof}

\sourcecorrection{OLAX-006}{મૂળની
\texttt{deduction-theorem-quantifiers.tex}, પંક્તિઓ 43--44માં
$((!A\land !C)\lif\lforall[x][!D(x)])\lif
(!A\lif(!C\lif\lforall[x][!D(x)]))$નું છેલ્લું જમણું કૌંસ ખૂટે છે.
અહીં તે ઉમેર્યું છે.}

\sourcecorrection{OLAX-007}{મૂળની
\texttt{deduction-theorem-quantifiers.tex}, પંક્તિ~48માં અંતિમ પરિણામ
$\Gamma\Proves !B$ લખાયું છે. પરંતુ અનુમાનના પગલાએ હમણાં જ
$\Gamma\Proves !A\lif(!C\lif\lforall[x][!D(x)])$ સ્થાપ્યું છે અને
$!B\ident !C\lif\lforall[x][!D(x)]$ છે; તેથી જરૂરી નિષ્કર્ષ
$\Gamma\Proves !A\lif !B$ લખ્યો છે.}

\begin{prob}
  \olref[fol][axd][ddq]{thm:deduction-thm-q}ની સાબિતી પૂરી કરો.
\end{prob}

\end{document}
